\documentclass[9pt]{article}
\usepackage[margin=0.66in]{geometry}
\usepackage{amsmath,amssymb,booktabs,microtype,xcolor}
\usepackage[T1]{fontenc}
\usepackage{lmodern}
\usepackage[colorlinks=true,allcolors=blue!55!black]{hyperref}
\setlength{\parindent}{0pt}
\setlength{\parskip}{3pt}
\title{The Work Cost of Reducing Dependence: A Harmonic-Oscillator Model}
\author{J. R. Landers}
\date{September 2026}
\begin{document}
\maketitle
\vspace{-1em}
\begin{abstract}
Can the energy required to reduce statistical dependence be calculated and tested? We define dependence using $r=p_{XY}/(p_Xp_Y)$ and mutual information, derive an exact reversible-work relation for a coupled harmonic model, and apply an explicitly approximate parameter bridge to a published nanomechanical experiment. Reported coupling settings then yield toy estimates of correlation, information reduction, and work, alongside a frequency check. This gives a concrete route from a general question to a measurable experiment.
\end{abstract}

\section{Dependence and a solvable equilibrium model}
Let $X,Y$ have joint density $p(x,y)$ and marginals $p_X,p_Y$. Their pointwise dependence ratio is
\begin{equation}
 r(x,y)=\frac{p(x,y)}{p_X(x)p_Y(y)},\qquad
 I(X;Y)=\iint p(x,y)\log r(x,y)\,dx\,dy .
\end{equation}
Independence is equivalent to $r=1$ almost everywhere and $I=0$. Since pointwise values of $r$ need not all move in the same direction, mutual information is a useful scalar measure of overall dependence.\cite{parrondo2015,sagawa2010} For a controlled Hamiltonian $H_\lambda(x,y)$ in a bath at temperature $T$, let $p_\lambda=Z_\lambda^{-1}e^{-H_\lambda/(k_BT)}$. The quasistatic work along a control path is
\begin{equation}
 W_{\rm rev}=\int_{\lambda_0}^{\lambda_1}\left\langle\partial_\lambda H_\lambda\right\rangle_\lambda d\lambda=F_{\lambda_1}-F_{\lambda_0},\qquad F_\lambda=-k_BT\log Z_\lambda.
\end{equation}
This makes the general calculation precise: predict $I_\lambda$ and $W_\lambda$ along a specified physical control path, then test both against data.\cite{jarzynski1997} There is no single energy number determined by $r$ alone; the Hamiltonian, bath, and path matter.

Consider two classical coordinates with potential energy
\begin{equation}
 E(x,y)=\frac{k}{2}(x^2+y^2)+gxy,\qquad k>0,\quad |g|<k.
\end{equation}
At equilibrium in a heat bath at temperature $T$, $p\propto e^{-E/(k_BT)}$. The covariance matrix is $k_BT\begin{pmatrix}k&g\\g&k\end{pmatrix}^{-1}$, so the correlation coefficient is $\rho=-g/k$. As this distribution is bivariate Gaussian,
\begin{equation}
 I=-\frac12\log(1-\rho^2)=-\frac12\log\!\left(1-\frac{g^2}{k^2}\right).
\end{equation}
Thus decreasing $|g|/k$ decreases dependence while leaving the marginal variances equal (their values change with $g$). The equilibrium free energy is $F(g)=C+(k_BT/2)\log(k^2-g^2)$, where $C$ is independent of $g$. For a quasistatic isothermal change $g_0\to g_1$,
\begin{equation}
 W_{\rm rev}=\Delta F=\frac{k_BT}{2}\log\!\left(\frac{k^2-g_1^2}{k^2-g_0^2}\right).
\end{equation}
For this particular family, $F(g)=C-k_BT I(g)$, hence the especially simple result
\begin{equation}
 W_{\rm rev}=k_BT\bigl[I(g_0)-I(g_1)\bigr].
\end{equation}
It says that reducing mutual information by one nat costs $k_BT$ of reversible work in this model. This is a model-specific law, not a universal identity for every system. Finite-rate dissipation and actuator costs add to the reversible work; changing temperature alone at fixed $g,k$ leaves $I$ unchanged.

\section{Connection to a measured optomechanical system}
A published experiment uses two spatially separated SiN membrane modes, each near $\omega_0\simeq2\pi\times400$ kHz. A cavity-mediated interaction tunes their effective coupling $\Lambda$. The article reports damping rates $\gamma_1=2\pi\times12$ Hz and $\gamma_2=2\pi\times6$ Hz, and illustrates coupling settings $\Lambda/2\pi=6,8,58$ Hz and approximately $85$ Hz. It reports coherent heat-flux oscillations in the strong-coupling regime and states that their frequency coincides with $2\Lambda$; oscillations disappear below the stated coupling threshold.\cite{yang2020,xuereb2015}

Using the reported settings, the direct numerical prediction is
\begin{center}
\begin{tabular}{rrr}
\toprule
$\Lambda/2\pi$ (Hz) & $2\Lambda/2\pi$ (Hz) & period (ms)\\
\midrule
6 & 12 & 83.3\\
8 & 16 & 62.5\\
58 & 116 & 8.62\\
85 & 170 & 5.88\\
\bottomrule
\end{tabular}
\end{center}
These frequencies follow algebraically from the reported coupling values and the paper's $2\Lambda$ relation. The article says the measured oscillation frequency coincides with this prediction, but does not give a table of numerical frequency estimates, so a residual or percent error cannot be calculated from its text. The damping rates correspond to amplitude-decay times $1/\gamma_1=13.3$ ms and $1/\gamma_2=26.5$ ms. At $58$ Hz coupling, the predicted flux period (8.62 ms) is shorter than these decay times, consistent with the reported coherent oscillations.

\paragraph{Putting numbers through the theory.} To connect the equilibrium theory to this device, make the explicit order-of-magnitude identification
\begin{equation}
 q:=\frac{|g|}{k}\approx\frac{\Lambda}{\omega_0}.
\end{equation}
This treats the dimensionless cross-stiffness as the ratio of coupling rate to resonant frequency. It is a modeling bridge, not a parameter identity established by the experiment. With $\omega_0/2\pi=400$ kHz, the reported settings give the following toy-model estimates:
\begin{center}
\begin{tabular}{rrrr}
\toprule
$\Lambda/2\pi$ (Hz) & $q$ & $|\rho|=q$ & $I$ (nats)\\
\midrule
6  & $1.50\times10^{-5}$ & $1.50\times10^{-5}$ & $1.13\times10^{-10}$\\
8  & $2.00\times10^{-5}$ & $2.00\times10^{-5}$ & $2.00\times10^{-10}$\\
58 & $1.45\times10^{-4}$ & $1.45\times10^{-4}$ & $1.05\times10^{-8}$\\
85 & $2.13\times10^{-4}$ & $2.13\times10^{-4}$ & $2.26\times10^{-8}$\\
\bottomrule
\end{tabular}
\end{center}
Here $I=-\frac12\log(1-q^2)\simeq q^2/2$. In this Gaussian analogue, reducing coupling from 85 to 58 Hz lowers estimated mutual information by about 53.5\%. The reversible-work formula can be evaluated in the same toy model: for 85$\to$58 Hz,
\begin{equation}
 W_{\rm rev}=\frac{k_BT}{2}\log\!\left(\frac{1-q_{58}^2}{1-q_{85}^2}\right)\simeq1.21\times10^{-8}k_BT\simeq5.0\times10^{-29}\;\mathrm{J}\quad(T=300\,\mathrm{K}).
\end{equation}
The work also follows directly from the model law: $k_BT\Delta I=1.21\times10^{-8}k_BT$. The resulting $5.0\times10^{-29}$ J is the toy model's ideal quasistatic cost for the estimated information reduction. It is not the optical energy used by the apparatus; that comparison is exactly what a calibrated experiment can test.

\section{Remarks}
The calculation now has a clear chain: infer a dependence statistic from a controlled coupling, use the model's $W_{\rm rev}=k_BT\Delta I$ relation, and compare both with measurements.\cite{sagawa2010,parrondo2015} The oscillator experiment is an exciting candidate because it tunes coupling and observes the response. Our estimate $|g|/k\approx\Lambda/\omega_0$ is a convenient first bridge, not an experimentally established identity.

The experiment drives the two membranes with different thermal reservoirs and tunes coupling using the cavity field. Its article reports mode frequencies, damping and coupling settings, heat-flux behavior, and effective-temperature trends; it says supporting data are available from the authors on request. To test the dependence-reduction theory, one needs synchronized displacement records (or numerical covariance matrices) at multiple coupling settings, plus calibration of the control power. Those data would permit estimating $p(x,y)$, $r(x,y)$, and $I$ with uncertainty, and comparing dependence before and after changing coupling. A quantitative control-energy estimate would additionally require the laser-power-to-$\Lambda$ calibration and the work/dissipation protocol. That gives us a concrete next experiment: estimate $p(x,y)$ and $I$ from the two motion traces as the laser tunes $\Lambda$, then compare the observed trend with the Gaussian prediction. In parallel, calibrated optical input and dissipation would let us put a real control-energy number beside the ideal-work benchmark. The oscillator platform already supplies the ingredients for a lively test of how coupling, dependence, and energy flow fit together.
\begin{thebibliography}{9}
\bibitem{jarzynski1997} C. Jarzynski, ``Nonequilibrium equality for free energy differences,'' \emph{Physical Review Letters} 78, 2690--2693 (1997). \href{https://doi.org/10.1103/PhysRevLett.78.2690}{doi:10.1103/PhysRevLett.78.2690}.
\bibitem{sagawa2010} T. Sagawa and M. Ueda, ``Generalized Jarzynski equality under nonequilibrium feedback control,'' \emph{Physical Review Letters} 104, 090602 (2010). \href{https://doi.org/10.1103/PhysRevLett.104.090602}{doi:10.1103/PhysRevLett.104.090602}.
\bibitem{parrondo2015} J. M. R. Parrondo, J. M. Horowitz, and T. Sagawa, ``Thermodynamics of information,'' \emph{Nature Physics} 11, 131--139 (2015). \href{https://doi.org/10.1038/nphys3230}{doi:10.1038/nphys3230}.
\bibitem{xuereb2015} A. Xuereb, A. Imparato, and A. Dantan, ``Heat transport in harmonic oscillator systems with thermal baths: application to optomechanical arrays,'' \emph{New Journal of Physics} 17, 055013 (2015). \href{https://doi.org/10.1088/1367-2630/17/5/055013}{doi:10.1088/1367-2630/17/5/055013}.
\bibitem{yang2020} C. Yang et al., ``Phonon heat transport in cavity-mediated optomechanical nanoresonators,'' \emph{Nature Communications} 11, 5933 (2020). \href{https://doi.org/10.1038/s41467-020-18426-4}{doi:10.1038/s41467-020-18426-4}.
\end{thebibliography}
\end{document}
